\Cref{tab:3-data} lists the six datasets. The main cross-system comparison
uses the first four; budget, lookup-table and scale tests cover all six.
Vogue contains fashion-image embeddings and arXiv contains abstract
embeddings. The two OpenAI datasets are different-dimensional Qdrant
DBpedia exports \cite{data_vogue,data_arxiv,data_openai1536,data_openai3072};
they share a corpus source. Alignment checks are reported in
\cref{sec:7-precision,sec:7-2}. BGE uses Italian Wikipedia, and Stella uses
the TREC24 BioGen corpus \cite{data_bge,data_stella}. The sizes below
describe the prepared exports. JHQ uses $M=d/8$, $B=B_r=8$, eight coarse-training iterations and
100,000 database rows for residual training. Cached runs reuse models
but re-encode the database; cold runs refit them.

\par\medskip\noindent\begin{minipage}{\linewidth}
\centering\captionsetup{hypcap=false}
\captionof{table}[Dataset and training configuration]{Data and JHQ training configuration. $n_q$ is the actual query count;
$n_t=39n_{\rm list}$ is the requested coarse-training sample size. The main comparison uses
the first four rows. All rows use squared L2 and $k=10$.}
\label{tab:3-data}
\small\setlength{\tabcolsep}{4pt}
% Generated by scripts/audit_ch3.py from the budget shape logs.
\begin{tabular*}{\linewidth}{@{\extracolsep{\fill}}lrrrrr@{}}
\toprule
Dataset & $N$ & $d$ & $n_q$ & $n_{\rm list}$ & $n_t$ \\
\midrule
Vogue & 932,328 & 768 & 1,000 & 4,096 & 159,744 \\
arXiv & 2,253,000 & 768 & 1,000 & 8,192 & 319,488 \\
OpenAI-1536 & 999,000 & 1,536 & 1,000 & 8,192 & 319,488 \\
OpenAI-3072 & 999,000 & 3,072 & 1,000 & 8,192 & 319,488 \\
BGE-M3 & 10,091,524 & 1,024 & 1,000 & 32,768 & 1,277,952 \\
Stella & 17,776,615 & 1,024 & 999 & 32,768 & 1,277,952 \\
\bottomrule
\end{tabular*}

\end{minipage}\par\medskip

OpenAI/BGE/Stella preparation removes vectors containing NaN or infinite
values (non-finite rows) and normalises nonzero vectors:
\[
 x\leftarrow x/\|x\|\quad(\|x\|>10^{-12}).
\]
Preparation scripts sample query rows without replacement using seed 42
and remove those rows from the base set. Stella retains 999 such corpus
queries, rather than the upstream 65-query set. This evaluates vector
neighbours of held-out corpus embeddings rather than relevance-labelled
user queries.

All methods use the same exported query and GT files. Preparation code
provides exhaustive squared-L2 GT over FP32 vectors through Faiss FlatL2
or a memory-bounded tiled scan. For query set $\mathcal Q$, returned ID
sets $\mathcal R_k(q)$ and stored GT prefixes $\mathcal T_k(q)$, recall is
\begin{equation}
\Recall@k=\frac{1}{|\mathcal Q|k}
\sum_{q\in\mathcal Q}|\mathcal R_k(q)\cap\mathcal T_k(q)|.
\label{eq:3-recall}
\end{equation}
GT files contain 100 neighbours for Vogue/arXiv and 20 for the others;
only the first ten count. Equal-distance alternatives outside that prefix
receive no automatic credit, and duplicate returned IDs add no hits.

Chapter 6 chooses a refinement budget from calibration queries without GT.
Even-indexed queries form the calibration pool and odd-indexed queries
are held out for testing, with zero-based indices:
\[
 \begin{aligned}
 \mathcal Q_{\rm cal}&=\{q_i:i\bmod2=0\},&
 \mathcal Q_{\rm hold}&=\{q_i:i\bmod2=1\},\\
 (|\mathcal Q_{\rm cal}|,|\mathcal Q_{\rm hold}|)
 &=\begin{cases}(500,499),&\text{Stella},\\(500,500),&\text{otherwise}.\end{cases}
 \end{aligned}
 \]
To examine how sampling changes the chosen budget, 64 random subsets
are drawn at each sample size $S\in\{32,64,128\}$ using seed 12345;
full-pool calibration uses all 500 calibration queries. Budget
selection compares returned IDs. GT is used afterwards to measure recall
and to construct a separate label-based budget comparison
(\cref{sec:6-2,sec:a-budget-contract}). Since the subsets reuse the same
query pool, these trials describe variation within that pool; Chapter 6
examines the assumptions needed to extend the result to future queries.
